% Lösungen Einheit 1 von 3 – Zählprinzip und Anordnungen (zusammenbau v0.1)
\einheitenkopf{Einheit 1 von 3 · Zählprinzip und Anordnungen}

\erg{8}{a) ja – Klassensprecher und Stellvertreter sind verschiedene Ämter \quad b) $4 \cdot 7 = 28$ \quad c) $160 : (5 \cdot 8) = 4$ \quad d) $5! = 5 \cdot 4 \cdot 3 \cdot 2 \cdot 1 = 120$ \quad e) $\frac{6!}{3! \cdot 2!} = 60$ (drei A und zwei N sind gleich) \quad f) $10^4 = 10\,000$ PINs; Anteil $\frac{5^4}{10^4} = 0{,}0625$ \quad g) $8 \cdot 7 \cdot 6 = 336$ \quad h) I $5! = 120$; II $5^5 = 3\,125$ (jeder Button geht an eines von fünf Kindern); III $\frac{8!}{3!} = 6\,720$}
\erg{9}{a) Fehler: Mia rechnet mit der Fakultät, obwohl ein Kind mehrere Aufkleber bekommen darf; jeder Aufkleber hat drei Empfänger: $3^3 = 27$ \quad b) Für den ersten Platz stehen alle Dinge zur Wahl, für jeden weiteren eines weniger, weil jedes nur einmal vorkommt; jede Reihenfolge zählt, also wird das fallende Produkt bis eins gebildet. \quad c) $6^4 = 1\,296$ Kombinationen; $1\,296 \cdot 3 = 3\,888$ s $\approx 64{,}8$ min}
